\chapter{Diffusion Models: from Thermodynamics to Scores}\label{ch:diffusion}

This chapter assembles the material of
Chapters~\ref{ch:statmech}--\ref{ch:ebm} into the modern theory of
generative diffusion. We present the discrete-time construction (DDPM), the
score-matching principle that makes training tractable, the continuous-time
SDE formulation that unifies both, and the identity --- Tweedie's --- that
converts the score into a Bayesian posterior expectation. The chapter closes
with the statistical-physics analyses of diffusion that motivate this
thesis, and states the gap our research chapters fill.

\section{The generative problem and the diffusion idea}

Given samples of an unknown distribution $p_0$ on $\R^K$, a generative model
must produce new samples. Diffusion models
\citep{sohl2015deep, ho2020denoising, song2021scorebased} solve this with a
two-process design inspired by non-equilibrium thermodynamics:
a fixed \emph{forward process} gradually corrupts data into pure noise, and
a learned \emph{reverse process} undoes the corruption step by step. The
forward process used throughout this thesis is the OU channel of
Section~\ref{sec:ou}: at diffusion time $t$,
\begin{equation}
  x_t \;=\; a\, e^{-t} + \sqrt{\Deltat}\;\varepsilon,
  \qquad a \sim p_0,\quad \varepsilon \sim \Nd(0, I),
  \qquad \Deltat = 1 - e^{-2t},
  \label{eq:forward-marginal}
\end{equation}
so that the noised marginal is the convolution
\begin{equation}
  p_t(x) \;=\; \int p_0(a)\; \Nd\!\big(x;\, a e^{-t},\, \Deltat I\big)\,
  \dd a .
  \label{eq:pt-convolution}
\end{equation}
As $t$ grows, $p_t$ interpolates from the data distribution to
$\Nd(0, I)$. By Anderson's theorem (Section~\ref{sec:reversal}), the
corruption is undone by the reverse SDE \eqref{eq:reverse-sde}, whose only
unknown ingredient is the score $\nabla_x \log p_t(x)$ of the noised
marginals. Generative diffusion is therefore, at bottom, \emph{score
estimation}: whoever knows $\nabla \log p_t$ for all $t$ can sample from
$p_0$.

\section{Denoising diffusion probabilistic models}\label{sec:ddpm}

The discrete-time formulation \citep{sohl2015deep, ho2020denoising}
discretises \eqref{eq:forward-marginal} into a Markov chain
$q(x_s | x_{s-1}) = \Nd(x_s; \sqrt{1 - \beta_s}\, x_{s-1}, \beta_s I)$ with
a variance schedule $\beta_1, \dots, \beta_S$, and learns a reverse chain
$p_\theta(x_{s-1} | x_s)$. Training maximises a variational lower bound on
the data log-likelihood --- the physicist will recognise the Gibbs
variational principle \eqref{eq:gibbs-variational} in modern dress.

\begin{toolbox}{The evidence lower bound of a diffusion model}
For any latent-variable model with joint $p_\theta(x_{0:S})$ and any
inference distribution $q(x_{1:S} | x_0)$, Jensen's inequality gives
\begin{align*}
  \log p_\theta(x_0)
  &= \log \int p_\theta(x_{0:S})\;\dd x_{1:S}
   = \log\, \E_{q}\!\left[
      \frac{p_\theta(x_{0:S})}{q(x_{1:S}|x_0)} \right]
  \;\geq\; \E_{q}\!\left[
      \log \frac{p_\theta(x_{0:S})}{q(x_{1:S}|x_0)} \right]
  =: \mathrm{ELBO}.
\end{align*}
Insert the two Markov factorisations
$p_\theta(x_{0:S}) = p(x_S) \prod_s p_\theta(x_{s-1}|x_s)$ and
$q(x_{1:S}|x_0) = \prod_s q(x_s|x_{s-1})$, then regroup telescopically using
Bayes' rule $q(x_s|x_{s-1}) = q(x_{s-1}|x_s, x_0)\, q(x_s|x_0) /
q(x_{s-1}|x_0)$:
\begin{equation*}
  -\mathrm{ELBO}
  = \underbrace{\kl{q(x_S|x_0)}{p(x_S)}}_{\text{prior term}}
  + \sum_{s=2}^{S}
    \underbrace{\E_q\, \kl{q(x_{s-1}|x_s, x_0)}{p_\theta(x_{s-1}|x_s)}}_{
      \text{one denoising step each}}
  \;-\; \E_q \log p_\theta(x_0 | x_1).
\end{equation*}
The decisive point: because the forward chain is Gaussian and conditioned on
\emph{both} endpoints, each target $q(x_{s-1} | x_s, x_0)$ is an explicit
Gaussian. Matching it with a Gaussian $p_\theta$ reduces each KL term to a
weighted regression of the posterior mean --- concretely, to predicting the
noise $\varepsilon$ that was added \citep{ho2020denoising}. The
thermodynamic reading of the ELBO gap as an entropy-production bound goes
back to \citet{sohl2015deep}. \hfill$\square$
\end{toolbox}

The practical upshot of \citet{ho2020denoising} is the loss
\begin{equation}
  \mathcal{L}(\theta)
  = \E_{s,\, a \sim p_0,\, \varepsilon}
    \big\| \varepsilon - \varepsilon_\theta\big(
      \underbrace{a\,e^{-t_s} + \sqrt{\Delta_{t_s}}\,\varepsilon}_{x_{t_s}},\; s
    \big) \big\|^2 ,
  \label{eq:ddpm-loss}
\end{equation}
a denoising regression. Its connection with the score is not incidental, as
the next section shows: predicting the noise \emph{is} estimating
$\nabla \log p_t$ up to scale.

\section{Score matching}\label{sec:score-matching}

Suppose we model the score directly by a function $s_\theta(x)$ and wish it
to match $\nabla \log p(x)$ for data from $p$. The naive objective
$\E_p \| s_\theta - \nabla \log p \|^2$ seems to require knowing the very
score we lack. Two classical results remove the obstruction.

\begin{toolbox}{Implicit score matching (Hyv\"arinen)}
Expand the square and keep only $\theta$-dependent terms:
\begin{equation*}
  \E_p \norm{s_\theta - \nabla \log p}^2
  = \E_p \norm{s_\theta}^2
  - 2\, \E_p \inner{s_\theta}{\nabla \log p}
  + \text{const}.
\end{equation*}
The cross term integrates by parts. Componentwise, using
$p\,\partial_i \log p = \partial_i p$ and assuming $p\, s_{\theta,i} \to 0$
at infinity:
\begin{equation*}
  \int p\; s_{\theta,i}\; \partial_i \log p \;\dd x
  = \int s_{\theta,i}\; \partial_i p \;\dd x
  = - \int p\; \partial_i s_{\theta,i} \;\dd x .
\end{equation*}
Therefore
\begin{equation*}
  \E_p \norm{s_\theta - \nabla \log p}^2
  = \E_p\Big[ \norm{s_\theta(x)}^2
    + 2\, \nabla \cdot s_\theta(x) \Big] + \text{const} :
\end{equation*}
an objective involving only the model and data samples
\citep{hyvarinen2005estimation}. Elegant, but the divergence term is
expensive for large networks --- which is why diffusion models use the
denoising variant instead. \hfill$\square$
\end{toolbox}

\begin{toolbox}{Denoising score matching (Vincent)}
Fix $t$ and consider the noised marginal \eqref{eq:pt-convolution}, written
as $p_t(x) = \int p_0(a)\, q_t(x|a)\,\dd a$ with Gaussian kernel
$q_t(x|a) = \Nd(x; a e^{-t}, \Deltat I)$. Then
\begin{equation*}
  \nabla_x \log p_t(x)
  = \frac{\int p_0(a)\, \nabla_x q_t(x|a)\, \dd a}{p_t(x)}
  = \int \underbrace{\frac{p_0(a)\, q_t(x|a)}{p_t(x)}}_{=\,p(a|x)}\;
    \nabla_x \log q_t(x|a)\; \dd a
  = \E\big[ \nabla_x \log q_t(x|a) \,\big|\, x \big].
\end{equation*}
The score of the \emph{marginal} is the posterior average of the score of
the \emph{kernel} --- and the latter is trivial:
$\nabla_x \log q_t(x|a) = -(x - a e^{-t})/\Deltat$. A standard
$L^2$-projection argument turns this into a regression: for any
$s_\theta$,
\begin{equation*}
  \E_{x \sim p_t} \norm{ s_\theta(x) - \nabla \log p_t(x) }^2
  = \E_{a, x} \norm{ s_\theta(x) - \nabla_x \log q_t(x|a) }^2
  + \text{const},
\end{equation*}
because conditional expectation is exactly the $L^2$ projection onto
functions of $x$ \citep{vincent2011connection}. Substituting the Gaussian
kernel score and $x = a e^{-t} + \sqrt{\Deltat}\varepsilon$ gives
$\nabla_x \log q_t = -\varepsilon / \sqrt{\Deltat}$: \emph{learning the
score of noisy data is learning to predict the noise}, and the DDPM loss
\eqref{eq:ddpm-loss} is denoising score matching in disguise.
\hfill$\square$
\end{toolbox}

The continuous-time synthesis \citep{song2019generative, song2021scorebased}
trains one network $s_\theta(x, t)$ against denoising score matching at all
noise levels and samples with the reverse SDE \eqref{eq:reverse-sde} or its
deterministic twin, the probability-flow ODE; an accessible exposition is
\citet{song2021blog}, and \citet{karras2022elucidating} systematise the
design space.

\section{Tweedie's identity: the score is a denoiser}\label{sec:tweedie}

The middle line of Vincent's derivation deserves its own section, because it
is the bridge on which the whole thesis walks. Rearranging
$\nabla \log p_t(x) = \E[-(x - a e^{-t})/\Deltat \mid x]$:
\begin{equation}
  \boxed{\;
  \nabla_x \log p_t(x)
  \;=\; \frac{e^{-t}\,\E[\,a \,|\, x\,] - x}{\Deltat}
  \quad\Longleftrightarrow\quad
  \E[\,a \,|\, x\,]
  = e^{t} \Big( x + \Deltat\, \nabla_x \log p_t(x) \Big).
  \;}
  \label{eq:tweedie}
\end{equation}
This is the Tweedie--Miyasawa formula of empirical Bayes
\citep{robbins1956empirical, miyasawa1961empirical, efron2011tweedie}:
the score of the noisy marginal and the Bayes-optimal denoiser
$\E[a|x]$ determine each other exactly, with no assumption whatsoever on
$p_0$. Three readings will recur:
\begin{enumerate}
  \item \emph{Geometric}: the score points from $x$ towards (the contracted
  image of) the posterior mean of the clean data --- denoising direction and
  score direction coincide.
  \item \emph{Statistical}: estimating the score is solving a Bayesian
  inference problem --- compute a posterior expectation under the joint
  \eqref{eq:hmm-joint}. Everything Chapter~\ref{ch:ebm} said about
  inference on graphs now applies to scores.
  \item \emph{Algorithmic}: for sequence data with Markov prior, the
  posterior lives on a tree (Proposition~\ref{prop:chain-tree}), so the
  score is computable by message passing --- the programme of
  Chapter~\ref{ch:bp}.
\end{enumerate}

For \emph{vector-valued} data $a \in \R^K$ noised coordinate-wise,
\eqref{eq:tweedie} holds componentwise with the \emph{full} posterior:
\begin{equation}
  \big( \nabla_x \log p_t(x) \big)_k
  \;=\; \frac{e^{-t}\, \E[\,a_k \,|\, x_0, \dots, x_{K-1}\,] - x_k}{\Deltat}.
  \label{eq:tweedie-joint}
\end{equation}
The conditioning is on \emph{all} coordinates: even though the noise is
independent across frames, the posterior mean of frame $k$ --- and hence the
$k$-th score component --- depends on every observation, because the prior
couples the frames. Equation \eqref{eq:tweedie-joint} is the precise reason
the \emph{joint} score differs from the per-frame marginal score, a
distinction that is easy to blur and whose importance for structured data
was impressed on this project in supervision; Chapter~\ref{ch:gaussian}
quantifies exactly how much coupling the conditioning induces, as a function
of $t$.

\section{The statistical physics of generative diffusion}
\label{sec:statphys-diffusion}

A recent line of work analyses diffusion models with the tools of
Chapter~\ref{ch:statmech}, in the regime where the score is \emph{exact} ---
either computed from a solvable $p_0$ or from the empirical measure of $n$
training points. Three findings frame our work.

\paragraph{Regimes and transitions.} For high-dimensional solvable data
(Gaussian mixtures, spin models), the reverse dynamics crosses sharp
dynamical regimes \citep{biroli2023generative, biroli2024dynamical}: a
\emph{speciation} time at which the trajectory commits to one mode ---
a spontaneous symmetry breaking in the sense of Section~\ref{sec:ising}
\citep[see also][]{raya2023spontaneous} --- and, when the score is learned
from finitely many samples, a \emph{memorisation} (collapse) time below
which the dynamics condenses onto training points. The phase boundaries are
computed with free-energy methods, in direct analogy with the random energy
model.

\paragraph{Sampling as physics.} \citet{ghio2024sampling} map generative
samplers --- flows, diffusions, autoregressive networks --- onto
statistical-physics problems and delimit where each is efficient, importing
the algorithmic-hardness picture of spin glasses
\citep{zdeborova2016statistical}.

\paragraph{The exactly solvable programme.} The common method of these
works is to choose $p_0$ rich enough to be interesting and simple enough
that $\nabla \log p_t$ is available in closed form, then study the
\emph{exact} generation dynamics. What none of them vary is the
\emph{internal structure of a single data point}: their coordinates are
exchangeable or i.i.d.\ conditional on a mode; there is no ordering, no
dynamics \emph{inside} the sample.

\section{The gap, restated}\label{sec:gap}

Real sequential data --- video, audio, trajectories --- have exactly the
structure the solvable programme leaves out: an internal (frame) time with a
Markovian dependence along it. For such data the score
\eqref{eq:tweedie-joint} is a smoothing posterior on a chain, an object with
seventy years of inference theory behind it
(Chapters~\ref{ch:stochproc}--\ref{ch:ebm}). Yet the questions a diffusion
modeller asks about it --- how does the coupling structure of
$\nabla \log p_t$ deform with $t$? can it be computed \emph{locally}? what
survives a Gaussian parametrisation of the messages when the data are not
Gaussian? what does that imply for score architectures? --- are not answered
by the classical literature, which fixes the noise level, nor by the
diffusion literature, which ignores internal structure. Answering them, in
the fully solvable setting of the AR(1) chain and its non-Gaussian
perturbations, is the contribution of
Chapters~\ref{ch:gaussian}--\ref{ch:bp}.
